% docs/documento_oficial/anexos.tex
\section*{Anexos}
\addcontentsline{toc}{section}{Anexos}

% ─────────────────────────────────────────────────────────────────────────────
\subsection*{Anexo A. Estructura del dataset experimental (\texttt{csv\_datos\_1.1.csv})}

El archivo CSV contiene los registros experimentales del semillero de la UCC utilizados por DIAPCE. Cada fila corresponde a un ensayo de resistencia a compresión con los siguientes campos:

\begin{table}[H]
  \centering
  \caption{Campos del dataset experimental}
  \begin{tabular}{lll}
    \hline
    Campo & Tipo & Descripción \\
    \hline
    \texttt{temperatura}     & Entero   & Temperatura de curado (°C): 10, 25 o 32 \\
    \texttt{humedad}         & Entero   & Humedad relativa (\%): 20, 50 o 70 \\
    \texttt{relacion\_ac}    & Decimal  & Relación agua/cemento: 0{,}40; 0{,}45 o 0{,}50 \\
    \texttt{edad\_dias}      & Entero   & Edad del ensayo: 7, 14 o 28 días \\
    \texttt{resistencia\_mpa}& Decimal  & Resistencia medida (MPa) \\
    \texttt{codigo\_aditivo} & Texto    & Código del aditivo: P0, P1, P2, P3, PP1, PP2 o PP3 \\
    \hline
  \end{tabular}
\end{table}

% ─────────────────────────────────────────────────────────────────────────────
\subsection*{Anexo B. Esquema completo de la base de datos local (DBML)}

\begin{lstlisting}[style=dbmlstyle, caption=Esquema completo de la base de datos local en DBML.]
Table users {
  id integer [pk, increment]
  email string [unique, not null]
  password string [not null]
}

Table projects {
  id integer [pk, increment]
  user_id integer [not null, ref: > users.id]
  project_name string [not null]
  selected_date string
  selected_image_path string
  creator_name string
  resistance_target integer [not null]
  temperature integer [not null]
  humidity integer [not null]
  relacion_ac float [not null]
  work_type string
  aditivo_id integer [ref: > Aditivos.id]
  resistencia_predicha_7d float
  resistencia_predicha_14d float
  resistencia_predicha_28d float
  created_at string
}

Table TiposAditivo {
  id integer [pk, increment]
  nombre string [unique, not null]
}

Table Productos {
  id integer [pk, increment]
  nombre_producto string [unique, not null]
  marca string
}

Table Aditivos {
  id integer [pk, increment]
  codigo string [unique, not null]
  porcentaje_aplicado string [not null]
  tipo_aditivo_id integer [not null, ref: > TiposAditivo.id]
  producto_id integer [not null, ref: > Productos.id]
}

Table ResultadosConcreto {
  id integer [pk, increment]
  temperatura integer [not null]
  humedad integer [not null]
  relacion_ac float [not null]
  edad_dias integer [not null]
  resistencia_mpa float [not null]
  aditivo_id integer [not null, ref: > Aditivos.id]

  indexes {
    (temperatura, humedad, relacion_ac, aditivo_id, edad_dias)
      [name: 'idx_busqueda_resistencia']
  }
}

Table mixtures {
  id integer [pk, increment]
  project_id integer [ref: > projects.id]
  name string
  description string
  cemento_kg float
  agua_litros float
  arena_kg float
  grava_kg float
  aditivo_porcentaje float
  created_at string
}
\end{lstlisting}

% ─────────────────────────────────────────────────────────────────────────────
\subsection*{Anexo C. Script SQL de creación de tablas (PostgreSQL)}

\begin{lstlisting}[style=sqlstyle, caption=Script DDL para PostgreSQL.]
CREATE TABLE TiposAditivo (
    id SERIAL PRIMARY KEY,
    nombre VARCHAR(50) UNIQUE NOT NULL
);
CREATE TABLE Productos (
    id SERIAL PRIMARY KEY,
    nombre_producto VARCHAR(100) UNIQUE NOT NULL,
    marca VARCHAR(50)
);
CREATE TABLE Aditivos (
    id SERIAL PRIMARY KEY,
    codigo VARCHAR(10) UNIQUE NOT NULL,
    porcentaje_aplicado VARCHAR(10) NOT NULL,
    tipo_aditivo_id INT NOT NULL REFERENCES TiposAditivo(id),
    producto_id INT NOT NULL REFERENCES Productos(id)
);
CREATE TABLE ResultadosConcreto (
    id SERIAL PRIMARY KEY,
    temperatura INT NOT NULL,
    humedad INT NOT NULL,
    relacion_ac DECIMAL(3, 2) NOT NULL,
    edad_dias INT NOT NULL,
    resistencia_mpa DECIMAL(10, 2) NOT NULL,
    aditivo_id INT NOT NULL REFERENCES Aditivos(id)
);
\end{lstlisting}
